% BLOCK11_STAGE2_MANAGED
% Lecture 11: smart pointers and ownership-aware APIs
\section{Лекция 11. Умные указатели и владение}
\label{sec:lecture-11}

В предыдущих лекциях мы уже встречали raw pointers, lifetime, RAII и
полиморфные interfaces. Остался центральный design-вопрос: \emph{кто отвечает
за уничтожение объекта?} Сам тип \texttt{T*} на этот вопрос не отвечает.
Указатель может быть лишь временным способом обратиться к уже существующему
object, а может случайно оказаться единственным адресом dynamically allocated
resource. Во втором случае ошибка ownership быстро превращается в leak,
double delete или dangling pointer.

\begin{importantbox}[Причинная линия]
Raw pointer не кодирует владение $\rightarrow$ ownership нужно выразить в API
$\rightarrow$ \texttt{unique\_ptr} задаёт одного owner $\rightarrow$
\texttt{shared\_ptr} задаёт совместное владение $\rightarrow$
\texttt{weak\_ptr} наблюдает shared object, не продлевая его lifetime.
\end{importantbox}

Главная цель smart pointers --- не «автоматически сделать любой pointer
безопасным». Они делают \term{ownership} явной частью типа и используют RAII:
когда owning smart pointer заканчивает lifetime, он выполняет соответствующий
release. При этом обычные references и raw pointers остаются полезными для
\emph{non-owning access}.

\subsection{Сначала словарь: object, access, ownership}

Разделим три понятия.

\term{Object} имеет lifetime. Пока lifetime не начался или уже закончился,
обращаться к нему как к живому object нельзя.

\term{Access} означает возможность обратиться к object. Reference
\texttt{T\&}, \texttt{const T\&} или raw pointer \texttt{T*} могут предоставлять
такой доступ, ничего не говоря об обязанности уничтожить object.

\term{Ownership} --- ответственность за окончание lifetime ресурса в нужный
момент. Для dynamically allocated object это обычно означает, что некоторый
owner обязан обеспечить ровно один корректный release.

\begin{whybox}[Почему raw pointer недостаточен как ownership contract?]
По declaration \texttt{Widget* p} caller не видит, должен ли он вызвать
\texttt{delete p}, запрещено ли это, может ли \texttt{p} быть null и кто должен
жить дольше кого. Эти сведения приходится хранить в комментариях и соглашениях.
Smart pointer переносит важную часть этого contract в type system.
\end{whybox}

В современном интерфейсе полезна следующая исходная эвристика:
\texttt{T\&}/\texttt{const T\&}: объект обязан существовать; функция его не владеет; \texttt{T*}: non-owning access, если допустимо отсутствие объекта или pointer semantics действительно удобна; \texttt{std::unique\_ptr<T>}: один owner; \texttt{std::shared\_ptr<T>}: несколько owners действительно должны совместно продлевать lifetime; \texttt{std::weak\_ptr<T>}: наблюдение за объектом, которым владеют \texttt{shared\_ptr}, без продления lifetime.

Это design guideline, а не правило синтаксиса стандарта. Legacy API может
использовать raw owning pointers, но такой contract труднее читать и проверять.

\subsection{Проблема naked \texttt{new}/\texttt{delete}}

Учебный код легко написать так:

\begin{cppcode}
Product* make_product() {
    return new Product{/*...*/};
}

void command() {
    Product* product = make_product();
    use(*product);               // а если здесь throw?
    delete product;
}
\end{cppcode}

Здесь ownership существует только в голове программиста. Early return,
exception или новый путь control flow легко пропускает \texttt{delete}.
Передача \texttt{product} другой функции также не отвечает на вопрос,
передаётся ли ownership.

RAII-идея из лекции~\ref{sec:lecture-10}: ответственность за release нужно
поместить в destructor объекта-owner. Standard library уже предоставляет такие
owners.

\subsection{\texttt{std::unique\_ptr}: один владелец}

\texttt{std::unique\_ptr<T>} находится в \headerfile{memory}. Он владеет
одним object (или пуст) и при уничтожении owner освобождает принадлежащий ему
object подходящим deleter.

\begin{cppcode}
#include <memory>

auto product = std::make_unique<Product>();
product->run();
\end{cppcode}

\texttt{std::make\_unique<T>(args...)} создаёт \texttt{T} и возвращает
\texttt{unique\_ptr<T>}. В C++14 и новее это обычный предпочтительный способ
получить unique ownership: raw owning \texttt{new} не оказывается между
несколькими выражениями user code.

\begin{standardbox}[Что гарантируется семантически]
\texttt{unique\_ptr} movable, но не copyable owner. Два независимых
\texttt{unique\_ptr} не могут обычным копированием начать владеть одним и тем же
object. Это важная часть ownership contract.
\end{standardbox}

Попытка:

\begin{cppcode}
auto first = std::make_unique<int>(42);
auto second = first; // compile error: copy запрещена
\end{cppcode}

проверяется настоящим negative compile test в эксперименте.

\subsection{Доступ через \texttt{unique\_ptr}}

Owner ведёт себя pointer-like:
\texttt{*owner} --- object; \texttt{owner->member()} --- member access; \texttt{owner.get()} --- raw pointer без передачи ownership; \texttt{if (owner)} --- проверка, есть ли owned object; \texttt{owner.reset()} --- заменить/освободить owned object; \texttt{owner.release()} --- отдать raw pointer и отказаться от владения.

\begin{warningbox}[Осторожно с \texttt{release()}]
После \texttt{release()} smart pointer уже не отвечает за cleanup. Caller снова
получает manual ownership problem. Это низкоуровневая операция для interop и
специальных contracts, а не обычный способ «достать pointer».
Для временного non-owning access нужен \texttt{get()}.
\end{warningbox}

\subsection{Передача unique ownership}

Ownership можно передать другому owner. Для уже существующей named variable
обычно нужна операция \texttt{std::move}:

\begin{cppcode}
void consume(std::unique_ptr<Job> job);

auto job = std::make_unique<Job>();
consume(std::move(job));

if (!job) {
    // ownership ушло в consume
}
\end{cppcode}

Здесь \texttt{std::move} используется только как минимальный bridge: он
разрешает выбрать move operation у movable type. Полная модель lvalue/rvalue,
value categories и то, что именно делает \texttt{std::move}, --- тема следующего
блока. Сейчас достаточно увидеть ownership contract: после передачи old owner
больше не владеет объектом.

Функция, которая \emph{принимает ownership}, может принимать
\texttt{std::unique\_ptr<T>} by value. Функция, которая только использует объект,
обычно не должна требовать unique pointer:

\begin{cppcode}
void print_product(const Product& product); // borrow
void maybe_print(const Product* product);   // nullable borrow
\end{cppcode}

Так caller не обязан менять способ владения только ради обычного чтения.

\subsection{Возврат ownership из factory}

В блоке 08 Factory Method специально обходился без smart pointers: creator
хранил concrete product внутри себя и возвращал non-owning reference. Теперь
мы можем выразить другой contract: factory создаёт новый polymorphic object и
\emph{передаёт caller единственное владение}.

\begin{cppcode}
class Product {
public:
    virtual ~Product() = default;
    virtual std::string format() const = 0;
};

class Creator {
public:
    virtual ~Creator() = default;
    std::unique_ptr<Product> create() const {
        return make_product();
    }

protected:
    virtual std::unique_ptr<Product> make_product() const = 0;
};
\end{cppcode}

Concrete creator возвращает, например,
\texttt{std::make\_unique<TextProduct>()}. Преобразование
\texttt{unique\_ptr<Derived>} в совместимый \texttt{unique\_ptr<Base>}
поддерживает polymorphic ownership.

Virtual destructor base здесь снова принципиален. В итоге unique owner хранит
\texttt{Product*}, но уничтожить нужно полный dynamic object.

\begin{whybox}[Почему factory return type полезен]
Return type \texttt{std::unique\_ptr<Product>} сразу сообщает caller:
«ты получаешь новый object и теперь отвечаешь за его lifetime, но cleanup
автоматический». Return \texttt{Product*} этого не выражал бы.
\end{whybox}

\subsection{Rule of Zero возвращается}

Если class хранит resource через \texttt{unique\_ptr} или standard containers,
ему часто не требуется manual destructor. Это продолжает Rule of Zero из
лекции~\ref{sec:lecture-06}: resource-managing member сам выполняет cleanup.

Но есть важное следствие. Class с \texttt{unique\_ptr} member обычно сам
становится non-copyable, потому что unique ownership нельзя копировать
автоматически. Позже move semantics позволит таким объектам передавать
ownership как значение.

\subsection{Массивы и unique ownership}

Существует специализация \texttt{std::unique\_ptr<T[]>} для dynamically allocated
arrays, и \texttt{std::make\_unique<T[]>(n)} умеет её создавать.

Для обычной динамической последовательности чаще удобен
\texttt{std::vector<T>}. Он хранит size, поддерживает iteration и value semantics.
\texttt{unique\_ptr<T[]>} полезен, когда задача действительно требует
низкоуровневого owning array buffer.

Smart pointer не заменяет container abstraction.

\subsection{\texttt{std::shared\_ptr}: совместное владение}

Иногда одного owner недостаточно по смыслу. Представим audit session, которая
должна оставаться живой, пока ею независимо пользуются console view и statistics
view. Ни один из них не подходит на роль единственного owner.

\begin{cppcode}
auto session = std::make_shared<AuditSession>("block-11");

ConsoleView console{session};
StatisticsView stats{session};
session.reset(); // views всё ещё совместно владеют session
\end{cppcode}

Копирование \texttt{shared\_ptr} создаёт ещё одного owner того же object.
Object уничтожается, когда исчезает последний owning \texttt{shared\_ptr}.

\begin{implementationbox}[Обычная реализация]
Обычно существует control block со счётчиками strong/weak references и
информацией, необходимой для destruction. \texttt{shared\_ptr} обращается к нему
при copy/destruction. Standard определяет observable semantics, но не требует
конкретного layout control block или конкретной ABI-структуры.
\end{implementationbox}

\texttt{std::make\_shared<T>} обычно предпочтительнее отдельного
\texttt{shared\_ptr<T>} с ручным \texttt{new T(...)}:
код проще, exception-safe, а обычная
implementation может объединить allocation object и control block. Последнее
--- optimization possibility, а не переносимая гарантия числа allocations для
всех реализаций.

\subsection{\texttt{use\_count()} не проектирует ownership}

\texttt{shared\_ptr::use\_count()} показывает число owning references в момент
наблюдения и полезен в учебных diagnostics. Но application logic не должна
обычно строиться на предположении «если count равен 1, то...»: между потоками и
изменениями owners это хрупкое знание, а сам count не объясняет правильную
модель владения.

Сначала проектируют, \emph{почему owners несколько}; затем выбирают
\texttt{shared\_ptr}.

\subsection{Цена shared ownership}

Shared ownership удобнее, но не бесплатно:
control block занимает память; copy/reset owners изменяют reference count; typical thread-safe reference-count updates могут требовать atomic operations; lifetime становится менее локально очевидным; циклы strong ownership не разрываются автоматически.

Поэтому policy проекта остаётся простой: если естественный owner один,
предпочитаем \texttt{unique\_ptr}. \texttt{shared\_ptr} появляется только при
реальном shared lifetime requirement.

\subsection{Цикл \texttt{shared\_ptr}}

Рассмотрим два objects, каждый из которых владеет другим:

\begin{cppcode}
struct Node {
    std::shared_ptr<Node> peer;
};

auto a = std::make_shared<Node>();
auto b = std::make_shared<Node>();
a->peer = b;
b->peer = a;
\end{cppcode}

После уничтожения local owners \texttt{a} и \texttt{b} каждый Node всё ещё
имеет одного strong owner: соседний Node. Reference counts не становятся нулём,
destructors не вызываются, возникает resource leak.

Это не dangling-pointer UB. Наоборот, objects остаются логически «живыми»
из-за неправильной ownership graph.

\subsection{\texttt{std::weak\_ptr}: наблюдение без владения}

\texttt{std::weak\_ptr<T>} связан с object, которым владеют
\texttt{shared\_ptr}, но сам не увеличивает strong owner count.

Он нужен для non-owning edges в shared graph:
parent link в tree, где parent владеет children; cache observation; subscriber, который не должен продлевать lifetime publisher; back-reference, разрывающий ownership cycle.

У \texttt{weak\_ptr} нет обычного \texttt{operator->}: сначала нужно временно
получить owning \texttt{shared\_ptr} через \texttt{lock()}.

\begin{cppcode}
std::weak_ptr<Session> observed = session;

if (auto alive = observed.lock()) {
    alive->print();
} else {
    // object уже уничтожен
}
\end{cppcode}

\texttt{lock()} атомарно относительно shared ownership semantics получает
\texttt{shared\_ptr}, если object ещё существует, либо пустой
\texttt{shared\_ptr}. \texttt{expired()} сообщает, что owning object уже исчез,
но для непосредственного использования всё равно удобнее проверять result
\texttt{lock()}.

\subsection{Правильное направление ownership graph}

Для tree естественная модель:

\begin{cppcode}
class Node {
    std::vector<std::shared_ptr<Node>> children_; // owning edges
    std::weak_ptr<Node> parent_;                  // non-owning back edge
};
\end{cppcode}

Это не означает, что любой tree обязан использовать \texttt{shared\_ptr}.
Если весь tree имеет одного owner, \texttt{unique\_ptr} для children может быть
лучше. Пример показывает именно принцип: ownership direction должен следовать
lifetime responsibility, а обратная навигация не обязана владеть.

\subsection{Raw pointer после smart pointer остаётся нормальным}

Smart pointers не отменяют raw pointers. Ошибка --- считать любой
\texttt{T*} owning pointer без ясного contract.

\begin{cppcode}
void inspect(const Product& product);

auto owner = std::make_unique<TextProduct>();
inspect(*owner);              // borrow через reference

const Product* view = owner.get(); // non-owning pointer
if (view != nullptr) {
    inspect(*view);
}
\end{cppcode}

\texttt{view} не должен пережить \texttt{owner}. Smart pointer гарантирует
cleanup owned object, но не может автоматически сделать безопасными все
заимствованные pointers/references.

\begin{warningbox}[Smart pointer не лечит dangling borrow]
Если сохранить \texttt{owner.get()}, затем уничтожить owner и после этого
разыменовать raw pointer, получится прежняя lifetime error. Ownership стало
явнее, но lifetime non-owning views всё ещё нужно соблюдать.
\end{warningbox}

\subsection{Как читать parameter types}

Практическая таблица API:

\begin{center}
\begin{tabularx}{\textwidth}{L{0.28\textwidth}YY}
\toprule
Parameter & Обычно означает & Вопрос caller \\
\midrule
\texttt{const T\&} & read-only borrow & object точно существует? \\
\texttt{T\&} & mutable borrow & mutation разрешена? \\
\texttt{const T*}/\texttt{T*} & nullable/non-owning view & возможен null? \\
\texttt{unique\_ptr<T>} by value & transfer unique ownership & готов отдать owner? \\
\texttt{const unique\_ptr<T>\&} & обычно слишком привязанный API & почему нужен именно owner type? \\
\texttt{shared\_ptr<T>} by value & функция становится co-owner & shared lifetime действительно нужен? \\
\texttt{const shared\_ptr<T>\&} & доступ к handle без нового owner & нужен именно handle contract? \\
\texttt{weak\_ptr<T>} & non-owning shared observation & object может исчезнуть? \\
\bottomrule
\end{tabularx}
\end{center}

Нет запрета принимать smart pointer by reference. Но если функция лишь вызывает
методы \texttt{T}, parameter \texttt{T\&}/\texttt{const T\&} обычно слабее
связывает API с конкретной ownership strategy.

\subsection{Return types и ownership}

Return type особенно хорошо выражает создание:
\texttt{T} --- вернуть самостоятельное value; \texttt{unique\_ptr<T>} --- вернуть единственного owner polymorphic/dynamic object; \texttt{shared\_ptr<T>} --- вернуть shareable owner, если API действительно обещает shared lifetime; \texttt{T*}/\texttt{T\&} --- обычно вернуть non-owning view на object, lifetime которого обеспечивается где-то ещё.

Возвращать \texttt{shared\_ptr} «на всякий случай» плохо: caller уже не может
понять, почему ownership должна быть shared, а code получает лишний control
block/reference-counting contract.

\subsection{Exception safety и smart pointers}

Smart pointers являются RAII owners, поэтому хорошо сочетаются с exceptions.

\begin{cppcode}
auto item = std::make_unique<Item>();
prepare(*item);       // may throw
store(std::move(item));
\end{cppcode}

Если \texttt{prepare} бросает до transfer, local unique owner уничтожится при
stack unwinding и освободит Item. Manual \texttt{delete} branch не нужен.

\texttt{make\_unique}/\texttt{make\_shared} также уменьшают число мест, где
raw resource временно существует без owner.

\subsection{Ошибочные способы создать несколько owners}

Нельзя создавать два независимых smart owners из одного raw owning pointer:

\begin{cppcode}
Widget* raw = new Widget;
std::shared_ptr<Widget> a{raw};
std::shared_ptr<Widget> b{raw}; // два разных control blocks: опасно
\end{cppcode}

Это не «два shared owners», а два независимых ownership systems, каждый считает,
что обязан уничтожить один object. Итог способен привести к double delete.

Правильный shared owner копируют:

\begin{cppcode}
auto a = std::make_shared<Widget>();
auto b = a; // один shared ownership group
\end{cppcode}

А unique owner вообще не копируют.

\subsection{Нужно ли всегда хранить \texttt{shared\_ptr} в members?}

Нет. Member type должен выражать responsibility конкретного class.
Если class обязан продлевать lifetime collaborator, shared owner оправдан.
Если collaborator гарантированно живёт снаружи дольше class, reference/raw
non-owning pointer может быть точнее.

Но такой borrow contract должен быть простым и доказуемым. Smart pointers не
заменяют design lifetime graph.

\subsection{Custom deleter: зачем существует}

Иногда resource освобождается не через обычный \texttt{delete}: C API может
требовать \texttt{fclose}, OS API --- специальную release function.
\texttt{unique\_ptr} и \texttt{shared\_ptr} поддерживают custom deleters.

В этой лекции мы не строим большой wrapper вокруг C resource: RAII-механизм уже
был изучен в блоке 10. Важно только увидеть общую идею: smart pointer --- это
policy-owning RAII handle, и способ destruction может быть частью его type/state.

\subsection{Потокобезопасность: аккуратная граница обещаний}

Shared ownership имеет thread-safety guarantees для операций над разными
\texttt{shared\_ptr} objects, относящимися к одному control block, но это не
делает сам \texttt{T} thread-safe. Два threads, меняющие один \texttt{T} через
shared owners без собственной synchronization, всё ещё могут создать data race.

Полную concurrency theory здесь не вводим.

\subsection{Стоимость: что именно мы покупаем}

\texttt{unique\_ptr} обычно проектируется как очень лёгкий owner: для обычного
deleter обычная implementation хранит один pointer, а destruction сводится к
проверке/удалению. Но стандарт не обещает универсально
\texttt{sizeof(unique\_ptr<T>) == sizeof(T*)}: custom deleter и implementation
могут менять representation.

\texttt{shared\_ptr} обычно дороже:
pointer/handle state; control block; reference-count updates; дополнительная indirection; более сложный lifetime.

Это цена именно shared ownership, а не причина избегать его там, где несколько
owners необходимы.

\begin{performancebox}[Не измеряем бессмысленный microbenchmark]
Главный выбор здесь семантический: один owner или несколько. Если ownership
model неверна, выигрыш в несколько наносекунд не исправляет design. Benchmark
уместен только после выбора корректной модели и при доказанном hot path.
Поэтому в блоке \texttt{BENCHMARKS=N/A}.
\end{performancebox}

\subsection{Эксперимент 1: ownership\_factory\_lab}

Каталог:
\filepath{examples/11\_smart\_pointers/01\_ownership\_factory\_lab}.

Эксперимент возвращается к Factory Method из блока 08, но теперь factory
создаёт новый polymorphic Product и возвращает
\texttt{std::unique\_ptr<Product>}.

Проверяются четыре идеи:
1) concrete creators возвращают unique ownership через base interface; 2) virtual destructor уничтожает полный dynamic object; 3) ownership можно передать через \texttt{std::move}, после чего old owner пуст; 4) helper, который лишь читает Product, принимает reference/raw non-owning view, а не smart pointer.

Отдельный \filepath{copy\_unique\_fail.cpp} действительно не компилируется:
копирование \texttt{unique\_ptr} запрещено.

\begin{shellcode}
cd examples/11_smart_pointers/01_ownership_factory_lab
cmake -S . -B build -DCMAKE_BUILD_TYPE=Debug
cmake --build build
ctest --test-dir build --output-on-failure
./run_checks.sh
\end{shellcode}

PASS означает, что все runtime cases успешны, negative compile действительно
отклонён compiler, а ownership/destructor trace совпадает с contract.

\subsection{Эксперимент 2: shared\_weak\_graph\_lab}

Каталог:
\filepath{examples/11\_smart\_pointers/02\_shared\_weak\_graph\_lab}.

Здесь есть три разных lifetime graph:
1) несколько independent views совместно владеют одним AuditSession; 2) два strong back-links создают цикл, и после внешнего scope destructor count остаётся нулевым; 3) tree хранит children через \texttt{shared\_ptr}, а parent back-link через \texttt{weak\_ptr}; цикл исчезает, destructors вызываются.

Также проверяется \texttt{weak\_ptr::lock}: пока strong owner существует,
\texttt{lock()} успешен; после его уничтожения weak observer сообщает expiration.

\begin{shellcode}
cd examples/11_smart_pointers/02_shared_weak_graph_lab
cmake -S . -B build -DCMAKE_BUILD_TYPE=Debug
cmake --build build
ctest --test-dir build --output-on-failure
./build/shared_weak_graph_lab
\end{shellcode}

Shared-cycle case намеренно демонстрирует logical leak в отдельном process.
Он не используется как production design и не UB. Исправленный weak
graph завершается нормальным destruction trace.

\subsection{Сводка выбора}

\begin{center}
\begin{tabularx}{\textwidth}{L{0.23\textwidth}YYY}
\toprule
Ситуация & Инструмент & Что владеет & Главный риск \\
\midrule
Один owner & \texttt{unique\_ptr} & один handle & premature transfer \\
Несколько owners & \texttt{shared\_ptr} & ownership group & cycles/overuse \\
Наблюдение shared object & \texttt{weak\_ptr} & никто дополнительно & object может исчезнуть \\
Короткий обязательный borrow & \texttt{T\&} & никто & owner должен жить дольше \\
Nullable borrow & \texttt{T*} & никто & null/dangling \\
\bottomrule
\end{tabularx}
\end{center}

\begin{importantbox}[Главный вывод]
Smart pointer выбирают не по правилу «raw pointer плохой», а по вопросу
ownership. Unique ownership --- default для dynamically owned object.
Shared ownership требует реального основания. Borrowing API обычно выражается
references/raw pointers и не должен случайно продлевать lifetime.
\end{importantbox}

Следующий блок объяснит move semantics и value categories. Тогда станет понятно,
почему \texttt{unique\_ptr} можно передавать, но нельзя копировать, и что именно
означает уже использованный здесь \texttt{std::move}.
